\documentclass[11pt,a4paper]{article}
\usepackage[T1]{fontenc}
\usepackage[utf8]{inputenc}
\usepackage[czech]{babel}
\usepackage{lmodern,microtype}
\usepackage{amsmath,amssymb,mathtools}
\usepackage{geometry}
\geometry{left=19mm,right=19mm,top=21mm,bottom=21mm}
\usepackage[dvipsnames]{xcolor}
\usepackage{enumitem,tabularx,booktabs}
\usepackage[most]{tcolorbox}
\usepackage{fancyhdr,lastpage}
\usepackage[hidelinks]{hyperref}
\definecolor{deep}{HTML}{203955}
\definecolor{teal}{HTML}{176E76}
\definecolor{slate}{HTML}{506173}
\definecolor{pale}{HTML}{F3F7F9}
\setlength{\parindent}{0pt}
\setlength{\parskip}{5pt}
\setlist[enumerate]{leftmargin=6mm,itemsep=5pt,topsep=3pt}
\pagestyle{fancy}\fancyhf{}
\fancyhead[L]{\sffamily\footnotesize\color{slate} MAT · 4. ročník · kombinatorika}
\fancyhead[R]{\sffamily\footnotesize\color{slate} Hodina \LESSONNUM\, / 112}
\fancyfoot[L]{\sffamily\footnotesize\color{slate} Výukový podklad · samostatný dokument}
\fancyfoot[R]{\sffamily\footnotesize\color{slate} \thepage/\pageref{LastPage}}
\renewcommand{\headrulewidth}{0.35pt}
\renewcommand{\footrulewidth}{0pt}
\newcommand{\sect}[1]{\par\vspace{2mm}{\Large\sffamily\bfseries\color{deep} #1}\par\vspace{-0.5mm}\noindent\color{teal}\rule{\linewidth}{0.7pt}\color{black}\vspace{1mm}}
\tcbset{theobox/.style={enhanced,breakable,colback=pale,colframe=teal!45!white,boxrule=.6pt,arc=1mm,left=3mm,right=3mm,top=2mm,bottom=2mm}}
\newcommand{\LESSONNUM}{32}
\begin{document}

{\Huge\sffamily\bfseries\color{deep} Permutace bez opakování; úvod do variací}\par
{\sffamily\large\color{teal} Výukový podklad · 32. vyučovací hodina}\par
\textbf{Hodinová dotace:} 45 minut\quad \textbf{Ročník:} 4. SPŠ\quad \textbf{Rozpis:} TP 2.1\par
\begin{tcolorbox}[theobox]\textbf{Výukový cíl.} Odvodit \(n!\) pro pořadí všech různých prvků a odlišit ho od uspořádaného výběru jen části prvků.\end{tcolorbox}
\sect{Teorie a vysvětlení}
\textbf{Permutace bez opakování} je uspořádání všech \(n\) různých prvků; jejich počet \(P(n)=n!\). Odvození: na první místo \(n\) možností, na druhé \(n-1\), a tak dále až poslední \(1\).\par
\textbf{Variace bez opakování} je uspořádaný výběr \(k\) různých prvků z \(n\), kde \(0\le k\le n\). Jejich počet \(V(n,k)=n(n-1)\cdots(n-k+1)=\frac{n!}{(n-k)!}\).\par
\textbf{Pozor na omezení.} „Dva stojí vedle sebe“ často řešíme spojením do bloku. „Nesmějí stát vedle sebe“ často vyřešíme odečtením nevhodných pořadí.
\sect{Řešené příklady}
\textbf{Příklad 1.} Kolika způsoby se do řady postaví šest různých studentů?\par
\textit{Řešení.} Jde o permutaci všech šesti osob: \(P(6)=6!=720\).\par\smallskip
\textbf{Příklad 2.} V šestičlenné řadě mají Adam a Bára stát vedle sebe. Kolik je pořadí?\par
\textit{Řešení.} Považujeme oba za jeden blok. Se čtyřmi ostatními lidmi tvoří pět objektů, které lze seřadit \(5!\) způsoby. Uvnitř bloku jsou dvě možnosti \(AB,BA\). Celkem \(2\cdot5!=240\).\par\smallskip
\textbf{Příklad 3.} Stejní studenti, Adam a Bára nesmějí stát vedle sebe.\par
\textit{Řešení.} Od všech \(6!=720\) pořadí odečteme \(240\) s dvojicí vedle sebe: \(720-240=480\).\par\smallskip
\textbf{Příklad 4.} Osm týmů obsazuje první tři místa. Kolik různých výsledků?\par
\textit{Řešení.} Nevybíráme všech osm týmů, ale pořadí medailí rozhoduje. Variace bez opakování: \(V(8,3)=8\cdot7\cdot6=336=\frac{8!}{5!}\).\par\smallskip
\textbf{Příklad 5.} Ze sedmi číslic \(1,\ldots,7\) vytvoř sedmiciferné číslo zakončené sudou cifrou bez opakování.\par
\textit{Řešení.} Na posledním místě jsou \(3\) možnosti: \(2,4,6\). Zbývajících šest číslic lze uspořádat \(6!\) způsoby. Výsledek \(3\cdot6!=2160\).\par\smallskip
\clearpage\sect{Procvičování pro studenty}
\begin{enumerate}
\item Kolik pořadí má pět různých knih?
\item Kolika způsoby získají z deseti závodníků zlato a stříbro dva různí?
\item Z pěti osob postav řadu, v níž je Petra vždy první.
\item Ze sedmi lidí sestav trojici vedoucí, zástupce a pokladník.
\item Šest lidí do řady, dva konkrétní musejí stát vedle sebe.
\end{enumerate}
\sect{Řešení a zdůvodnění k procvičování}
\begin{enumerate}
\item \(5!=120\).
\item \(V(10,2)=10\cdot9=90\).
\item Petra pevně první, zbytek má \(4!=24\) pořadí.
\item Pořadí funkcí je důležité, tedy \(7\cdot6\cdot5=210\).
\item Jeden blok plus čtyři osoby: \(5!\cdot2!=240\).
\end{enumerate}
\sect{Metodické poznámky}
\textbf{Typické chyby:}
\begin{itemize}[leftmargin=5mm,itemsep=1pt]
\item Permutace používají všechny prvky, variace mohou jen část.
\item „Skupina“ a „přiřazení funkcí“ nejsou totéž.
\item Blok dvou lidí má zpravidla také vnitřní pořadí.
\end{itemize}
\textbf{Doporučený průběh 45 minut:} 0–10 odvození \(n!\); 10–23 podmínky/bloky; 23–33 variace bez opakování; 33–42 cvičení; 42–45 výstup.\par
\textbf{Návaznost:} Nyní dovolíme opakované nerozlišitelné prvky; obyčejné \(n!\) pak výsledky nadhodnocuje.\par
\vspace{2mm}\noindent{\footnotesize\color{slate} \textbf{Východiska:} rozpis výuky MAT, 4. ročník, hodiny 29--41; tematický plán SPŠ; srovnání tematického rozsahu a náročnosti s realisticky.cz, oddíl Kombinatorika (lekce 090101--090119). Výklad, příklady a řešení jsou nově autorsky sestavené, nikoli převzatá sbírka.}\par
\end{document}